% ============================================================
% The Collision Prime Number Theorem
% ============================================================
%
% Paper E: combines research notes 25 and 26.
% Target: arXiv submission.
%
% ============================================================

\documentclass[12pt]{amsart}

\usepackage{amsmath,amssymb,amsthm}
\usepackage{booktabs}
\usepackage{microtype}
\usepackage{url}
\usepackage{hyperref}

\newtheorem{theorem}{Theorem}[section]
\newtheorem{lemma}[theorem]{Lemma}
\newtheorem{proposition}[theorem]{Proposition}
\newtheorem{corollary}[theorem]{Corollary}
\theoremstyle{definition}
\newtheorem{definition}[theorem]{Definition}
\theoremstyle{remark}
\newtheorem{remark}[theorem]{Remark}

\title{The Collision Prime Number Theorem}

\author{Alexander S.\ Petty}

\email{alexander.petty@gmail.com}
\date{March 2026}
\subjclass[2020]{11A63, 11N13, 11M06, 11A07}

\keywords{Collision invariant, digit function, Dirichlet
$L$-functions, prime number theorem, Siegel zeros}

\begin{document}
\begin{abstract}
For a prime $p$ not dividing a positional base~$b$, the
\emph{collision count} $C(b \bmod p)$ records the number of
residues $r$ for which the digit function
$\delta(r) = \lfloor br/p \rfloor$ satisfies
$\delta(r) = \delta(br \bmod p)$. The \emph{signed collision
count} $S(p) = C(b \bmod p) - \lfloor(p{-}1)/b\rfloor$ is
proved to be an integer determined by $p \bmod b^2$, taking
values in $[-(b{-}1),\; b{-}2]$ with the antisymmetry
$S(a) + S(b^2{-}a) = -1$. The $\varphi(b^2)$ values form
the \emph{collision periodic table}, a finite invariant of
the base.

The centered collision deviation
$S^{\circ}(p)$ is a linear combination of primitive odd
Dirichlet characters modulo~$b^2$. Three consequences follow.
First, the \emph{collision prime number theorem}:
$\sum_{p \le x} S^{\circ}(p) =
O(x \exp(-c\sqrt{\log x}))$,
with an effective constant $c > 0$ for all $b \ge 3$
(no Siegel zero caveat). Second, the \emph{collision
variance}: $\sum |S^{\circ}(p)|^2 \sim \sigma^2 x / \log x$
for an explicit constant $\sigma^2$. Third, the periodic table encodes the values
$L(1, \chi)$ for all primitive odd characters
modulo~$b^2$: the classical formula
$|B_1(\overline{\chi})| = (b/\pi) |L(1,\chi)|$
converts the table's Fourier coefficients to $L$-values.
In the composite bases tested ($b = 6, 10, 12$), the
extraction has a single constant $1/(\pi \varphi(b))$.
\end{abstract}
\maketitle

% ============================================================
\section{Introduction}

The digit function $\delta(r) = \lfloor br/p \rfloor$ assigns
to each residue $r \in \{1, \ldots, p{-}1\}$ a digit
$d \in \{0, \ldots, b{-}1\}$. Multiplication by $g \bmod p$
permutes the residues. The \emph{collision count}
$C(g) = \#\{r : \delta(r) = \delta(gr \bmod p)\}$ measures
how many residues share a digit with their image under this
permutation.

The preceding papers in this series~\cite{paperA,paperB,paperC}
established that the collision count decomposes into Dirichlet
characters modulo~$b^2$. This paper proves three consequences
of that decomposition and exhibits the finite object that
encodes the collision structure of all primes.

The collision count at the specific multiplier $g = b$ has a
closed form. Write $p = b^2 Q_2 + a$ with
$0 < a < b^2$ and $\gcd(a, b) = 1$. Define
\begin{equation}\label{eq:f}
f(a) = \#\{d \in \{0, \ldots, b{-}1\} :
(b{+}1)da \bmod b^2 \ge b^2 - a\}.
\end{equation}

\begin{proposition}[Collision formula]\label{prop:formula}
For a prime $p > b^2$ with $p \nmid b$:
\[
C(b \bmod p) = b \cdot Q_2 + f(a) - 1,
\]
where $Q_2 = \lfloor p / b^2 \rfloor$ and $a = p \bmod b^2$.
\end{proposition}

The signed collision count
$S(p) = C(b \bmod p) - \lfloor(p{-}1)/b\rfloor$
satisfies $S = f(a) - 1 - \lfloor(a{-}1)/b\rfloor$, depending
only on $a = p \bmod b^2$.

% ============================================================
\section{The Periodic Table}

\begin{theorem}[Finite determination]\label{thm:fingerprint}
For any base $b \ge 2$ and prime $p > b^2$ with $p \nmid b$,
the signed collision count $S(p)$ is an integer determined by
$p \bmod b^2$.
\end{theorem}

\begin{proof}
$S(p) = f(a) - 1 - \lfloor(a{-}1)/b\rfloor$ by
Proposition~\ref{prop:formula}, where $a = p \bmod b^2$.
\end{proof}

\begin{definition}
The \emph{collision periodic table} is the function
$\mathcal{T} \colon (\mathbb{Z}/b^2\mathbb{Z})^{*}
\to \mathbb{Z}$ defined by $\mathcal{T}(a) = S(p)$ for
any prime $p \equiv a \pmod{b^2}$ with $p > b^2$.
\end{definition}

\begin{theorem}[Antisymmetry]\label{thm:antisym}
For every base $b \ge 3$ and every coprime $a \bmod b^2$:
\[
\mathcal{T}(a) + \mathcal{T}(b^2 - a) = -1.
\]
\end{theorem}

\begin{proof}
Three identities.

\emph{Complementary counting.}
Let $r_d = (b{+}1)da \bmod b^2$ for $d = 0, \ldots, b{-}1$.
For $d \ge 1$: $r_d \ne 0$ (since $\gcd(a(b{+}1), b^2) = 1$).
The complement $b^2 - r_d = (b{+}1)d(b^2{-}a) \bmod b^2$.
The condition $r_d \ge b^2 - a$ for $f(a)$ and the condition
$b^2 - r_d \ge a$ for $f(b^2{-}a)$ are complementary: every
$d \ge 1$ satisfies exactly one, except when $r_d = b^2 - a$
(both). Including $d = 0$ ($r_0 = 0$, neither):
$f(a) + f(b^2{-}a) = b$.

\emph{Floor identity.}
$(a{-}1) + (b^2{-}a{-}1) = b^2 - 2$, with remainders
$R + R' = b - 2 < b$ modulo~$b$, so
$\lfloor(a{-}1)/b\rfloor + \lfloor(b^2{-}a{-}1)/b\rfloor
= b - 1$.

\emph{Assembly.}
$\mathcal{T}(a) + \mathcal{T}(b^2{-}a) =
[f(a) + f(b^2{-}a)] - 2 -
[\lfloor(a{-}1)/b\rfloor + \lfloor(b^2{-}a{-}1)/b\rfloor]
= b - 2 - (b{-}1) = -1$.
\end{proof}

\begin{corollary}\label{cor:range}
$\mathcal{T}$ takes values in $[-(b{-}1),\; b{-}2]$.
Exactly $\varphi(b^2)/2$ classes are negative. The mean
of $\mathcal{T}$ is $-1/2$.
\end{corollary}

\begin{proof}
For the lower bound: $f(a) \ge 1$ because $d = b{-}1$
always contributes ($(b{+}1)(b{-}1) \equiv -1 \pmod{b^2}$,
so $r_{b-1} = b^2 - a \ge b^2 - a$). Thus
$\mathcal{T}(a) = f(a) - 1 - \lfloor(a{-}1)/b\rfloor
\ge 0 - (b{-}1) = -(b{-}1)$. By antisymmetry,
$\mathcal{T}(a) = -1 - \mathcal{T}(b^2{-}a)
\le -1 + (b{-}1) = b - 2$.

Each pair $(a, b^2{-}a)$ sums to $-1 < 0$, so at least
one member is negative. Since $\mathcal{T}(a) \ge 0$
implies $\mathcal{T}(b^2{-}a) \le -1 < 0$, every pair
contributes exactly one negative member. The mean follows:
$\varphi(b^2)/2$ pairs each summing to $-1$.
\end{proof}

\begin{table}[ht]
\centering
\begin{tabular}{r*{10}{r}}
\toprule
& $0$ & $1$ & $2$ & $3$ & $4$ & $5$ & $6$ & $7$ & $8$ & $9$ \\
\midrule
$0$ & $.$ & $0$  & $.$ & $+2$ & $.$ & $.$ & $.$ & $0$  & $.$ & $+8$ \\
$1$ & $.$ & $-1$ & $.$ & $-1$ & $.$ & $.$ & $.$ & $+1$ & $.$ & $-1$ \\
$2$ & $.$ & $0$  & $.$ & $-2$ & $.$ & $.$ & $.$ & $+6$ & $.$ & $0$ \\
$3$ & $.$ & $-1$ & $.$ & $-1$ & $.$ & $.$ & $.$ & $-3$ & $.$ & $-1$ \\
$4$ & $.$ & $-4$ & $.$ & $0$  & $.$ & $.$ & $.$ & $-2$ & $.$ & $0$ \\
$5$ & $.$ & $-1$ & $.$ & $+1$ & $.$ & $.$ & $.$ & $-1$ & $.$ & $+3$ \\
$6$ & $.$ & $0$  & $.$ & $+2$ & $.$ & $.$ & $.$ & $0$  & $.$ & $0$ \\
$7$ & $.$ & $-1$ & $.$ & $-7$ & $.$ & $.$ & $.$ & $+1$ & $.$ & $-1$ \\
$8$ & $.$ & $0$  & $.$ & $-2$ & $.$ & $.$ & $.$ & $0$  & $.$ & $0$ \\
$9$ & $.$ & $-9$ & $.$ & $-1$ & $.$ & $.$ & $.$ & $-3$ & $.$ & $-1$ \\
\bottomrule
\end{tabular}
\caption{The collision periodic table in base~$10$.
Row is the tens digit, column is the units digit of
$p \bmod 100$. Dots indicate non-coprime classes. For
any prime $p > 100$, the entry gives
$S = C(10 \bmod p) - \lfloor(p{-}1)/10\rfloor$.}
\label{tab:periodic}
\end{table}

% ============================================================
\section{The Collision Prime Number Theorem}

The decomposition theorem~\cite{paperB} gives
\begin{equation}\label{eq:decomp}
S^{\circ}(p) = \sum_{\chi} \widehat{c}(\chi)\, \chi(p),
\end{equation}
a sum over primitive odd characters modulo~$b^2$, where the
\emph{centered} deviation $S^{\circ} = S - \overline{S}_R$
removes the class mean $\overline{S}_R$ (which depends on
$R = (p{-}1) \bmod b$). The coefficients $\widehat{c}(\chi)$
are explicitly determined~\cite{paperC}.

\begin{theorem}[Collision PNT]\label{thm:pnt}
For any base $b \ge 2$:
\[
\sum_{\substack{p \le x \\ p \nmid b}} S^{\circ}(p)
= O\!\left(x \exp\!\left(-c\sqrt{\log x}\right)\right),
\]
where $c > 0$ depends on $b$ (effective for $b \ge 3$
by Theorem~\ref{thm:siegel}).
\end{theorem}

\begin{proof}
The sum equals
$\sum_{\chi} \widehat{c}(\chi) \sum_{p \le x} \chi(p)$.
The outer sum is finite ($\le \varphi(b^2)/2$ terms).
Each inner sum is
$O(x \exp(-c_1\sqrt{\log x}))$ by the classical
zero-free region~\cite{davenport}.
\end{proof}

\begin{theorem}[Siegel immunity]\label{thm:siegel}
For $b \ge 3$, no real primitive character modulo~$b^2$
exists. The constant $c$ in Theorem~\ref{thm:pnt} is
effective.
\end{theorem}

\begin{proof}
A quadratic character modulo $p^{2a}$ ($p$ odd) has
conductor~$p$, not~$p^{2a}$: it is the lift of the
Legendre symbol. At $p = 2$ with $2a \ge 4$, the
conductor is at most~$8$. Since $b^2$ is a perfect
square and $b \ge 3$, every quadratic character
modulo~$b^2$ is imprimitive. The decomposition
uses only primitive characters.
\end{proof}

\begin{theorem}[Collision variance]\label{thm:variance}
$\displaystyle
\sum_{\substack{p \le x \\ p \nmid b}} |S^{\circ}(p)|^2
= \sigma^2 \cdot \frac{x}{\log x}
+ O\!\left(\frac{x}{\log^2 x}\right),
$
where $\sigma^2 = \sum_{\chi} |\widehat{c}(\chi)|^2$.
\end{theorem}

\begin{proof}
The finitely many primes $p \le b^2$ contribute $O(1)$
and are absorbed into the error term. For $p > b^2$:
$S^{\circ}(p)$ depends only on $p \bmod b^2$
(Theorem~\ref{thm:fingerprint}), so the sum decomposes
over residue classes:
\[
\sum_{p \le x} |S^{\circ}(p)|^2 =
\sum_{\substack{a \bmod b^2 \\ \gcd(a,b) = 1}}
|S^{\circ}(a)|^2 \cdot \pi(x; b^2, a),
\]
where $\pi(x; q, a) = \#\{p \le x : p \equiv a \pmod{q}\}$.
By the PNT for arithmetic progressions with fixed
modulus~$b^2$:
$\pi(x; b^2, a) = x / (\varphi(b^2) \log x) +
O(x / \log^2 x)$ uniformly in~$a$. Summing:
\[
\sum_{p \le x} |S^{\circ}(p)|^2 =
\frac{x}{\varphi(b^2) \log x}
\sum_a |S^{\circ}(a)|^2 +
O\!\left(\frac{x}{\log^2 x}\right).
\]
By Parseval on $(\mathbb{Z}/b^2\mathbb{Z})^{*}$ (with
the normalization $\widehat{c}(\chi) =
\frac{1}{\varphi(b^2)} \sum_a S^{\circ}(a)\,
\overline{\chi(a)}$ from~\cite{paperB}):
$\frac{1}{\varphi(b^2)} \sum_a |S^{\circ}(a)|^2 =
\sum_{\chi} |\widehat{c}(\chi)|^2 = \sigma^2$.
\end{proof}

In base~$10$ at lag~$1$: $\sigma^2 \approx 7.26$, giving
typical fluctuations $|S^{\circ}| \approx 2.7$.

% ============================================================
\section{The $L$-value Extraction}

The collision periodic table encodes $L$-function values.
The extraction works in every base.

\begin{theorem}[$L$-value extraction]\label{thm:lvalue}
For every base $b \ge 2$, the values $L(1, \chi)$ for all
primitive odd characters modulo~$b^2$ are determined by
the collision periodic table via three steps:
\begin{enumerate}
\item Fourier inversion of $\mathcal{T}(a)$ on
      $(\mathbb{Z}/b^2\mathbb{Z})^{*}$ recovers
      $\widehat{c}(\chi)$. (The class mean
      $\overline{S}_R$ depends on $p \bmod b$, so
      primitive characters modulo~$b^2$ are orthogonal
      to it; inversion of $\mathcal{T}$ and of
      $S^{\circ}$ yield the same $\widehat{c}$.)
\item The classical formula
      $|B_1(\overline{\chi})| = (b / \pi) \cdot
      |L(1, \chi)|$ gives $|L(1, \chi)|$ from
      the (computable) generalized Bernoulli
      number~\cite{davenport}.
\item The coefficient formula~\cite{paperC}
      $\widehat{c}(\chi) = -B_1(\overline{\chi})
      \cdot \overline{S_G(\chi)} / \varphi(b^2)$
      provides internal consistency: the bin sum
      $S_G(\chi)$ is recovered as
      $S_G = -\varphi(b^2) \widehat{c} /
      B_1(\overline{\chi})$.
\end{enumerate}
The system is overdetermined ($\varphi(b^2)$ table entries,
$\varphi(b^2)/2$ character coefficients), providing
verification at each step.
\end{theorem}

In the composite bases tested, the extraction admits a
single proportionality constant:

\begin{remark}[Composite bases]\label{rem:composite}
Computation in bases $b \in \{6, 10, 12\}$ shows
$|S_G(\chi)| = |B_1(\overline{\chi})|$ for all
primitive odd~$\chi$, giving
\[
|\widehat{c}(\chi)| =
\frac{|B_1(\overline{\chi})| \cdot |L(1, \chi)|}
{\pi \cdot \varphi(b)}.
\]
In base~$10$: the constant is $1/(4\pi) = 0.07958$,
verified for all $8$ primitive odd characters
modulo~$100$ (Table~\ref{tab:lvalues}). The identity
$|S_G| = |B_1|$ for composite bases is observed, not
proved.
\end{remark}

For prime bases $b$ with $\varphi(b) \le 4$
(i.e., $b \in \{3, 5\}$), the ratio
$|S_G(\chi)| / |B_1(\overline{\chi})|$ is also constant
across characters ($b$ for $b = 3$, $\sqrt{b}$ for
$b = 5$). For prime bases with $\varphi(b) \ge 6$
(starting at $b = 7$), the ratio varies by character.
In all cases, the extraction
(Theorem~\ref{thm:lvalue}) is exact.

\begin{table}[ht]
\centering
\begin{tabular}{rrrrr}
\toprule
$k$ & $|\widehat{c}|$ & $|B_1|$ & $|L(1,\chi_k)|$ &
$\frac{|\widehat{c}|}{|B_1| \cdot |L|}$ \\
\midrule
$2$  & $0.462$ & $4.298$ & $1.350$ & $0.0796$ \\
$4$  & $0.362$ & $3.804$ & $1.195$ & $0.0796$ \\
$6$  & $0.238$ & $3.087$ & $0.970$ & $0.0796$ \\
$8$  & $0.138$ & $2.351$ & $0.739$ & $0.0796$ \\
$12$ & $0.138$ & $2.351$ & $0.739$ & $0.0796$ \\
$14$ & $0.238$ & $3.087$ & $0.970$ & $0.0796$ \\
$16$ & $0.362$ & $3.804$ & $1.195$ & $0.0796$ \\
$18$ & $0.462$ & $4.298$ & $1.350$ & $0.0796$ \\
\bottomrule
\end{tabular}
\caption{$L$-value extraction in base~$10$. The ratio
$|\widehat{c}| / (|B_1| \cdot |L|) = 1/(4\pi) = 0.07958$
for all $8$ primitive odd characters modulo~$100$.}
\label{tab:lvalues}
\end{table}

% ============================================================
\section{Remarks}

\subsection*{Collision count from the table}

The periodic table determines $C(b \bmod p)$ directly:
$C = \lfloor(p{-}1)/b\rfloor + \mathcal{T}(p \bmod b^2)$.
The table replaces an $O(p)$ computation with a lookup
in a fixed table of $\varphi(b^2)$ integers.

\subsection*{Universality across bases}

The antisymmetry $\mathcal{T}(a) + \mathcal{T}(b^2{-}a) = -1$,
the range $[-(b{-}1), b{-}2]$, and the negative majority
$\varphi(b^2)/2$ hold for every base $b \ge 3$. The table
and its structural properties are base-dependent in their
entries but base-independent in their form.

\subsection*{The sign balance}

The distribution of $S^{\circ}(p)$ across primes is discrete:
$S^{\circ}$ takes finitely many values, one per coprime class
modulo~$b^2$. The sign of $S^{\circ}$ is determined by
$p \bmod b^2$. In base~$10$, exactly $20$ classes are positive
and $20$ are negative.

\subsection*{Relation to the class number formula}

The $L$-value extraction parallels Dirichlet's class number
formula~\cite{davenport}: a finite arithmetic quantity (the
class number, or the periodic table) determines $L(1, \chi)$
through a formula involving~$\pi$. The class number formula
extracts $L$-values from ideal class structure. The collision
periodic table extracts them from digit collisions.

\subsection*{The $S_G$ landscape}

Computation in composite bases ($b = 6, 10, 12$) shows
$|S_G(\chi)| = |B_1(\overline{\chi})|$, giving a single
constant $1/(\pi \varphi(b))$.
For prime bases with $\varphi(b) \le 4$: the ratio
$|S_G|/|B_1|$ is constant ($b$ for $b = 3$,
$\sqrt{b}$ for $b = 5$). For prime bases with
$\varphi(b) \ge 6$: the ratio varies by character. The
transition at $\varphi(b) = 6$ reflects the point where the
character group becomes rich enough to resolve structural
differences between digit bins. The extraction is exact
in all cases; the single-constant form is a simplification
available for composite bases and small primes.

% ============================================================
\begin{thebibliography}{9}

\bibitem{paperA}
A.~S.~Petty,
The collision invariant,
preprint, 2026.
\href{https://arxiv.org/abs/2603.17408}{arXiv:2603.17408}.

\bibitem{paperB}
A.~S.~Petty,
The collision transform,
preprint, 2026.
\href{https://arxiv.org/abs/2603.17432}{arXiv:2603.17432}.

\bibitem{paperC}
A.~S.~Petty,
The collision spectrum,
preprint, 2026.

\bibitem{davenport}
H.~Davenport,
\emph{Multiplicative Number Theory},
3rd~ed., Springer, 2000.

\bibitem{iwaniec}
H.~Iwaniec and E.~Kowalski,
\emph{Analytic Number Theory},
Amer.\ Math.\ Soc., 2004.

\bibitem{nfield}
A.~S.~Petty,
\emph{nfield: A structural analysis engine for fractional
field invariants}, software.
\url{https://github.com/alexspetty/nfield}

\end{thebibliography}

\end{document}
